\SGASubsection{4.4}{Osculating position}
\SGARefTarget{sga3:xxvi:subsection:4:4}

\begin{sgastatement}{\SGAInlineRefTarget{sga3:xxvi:proposition:4-4-1}Proposition 4.4.1}
Let \(S\) be a scheme, \(G\) an \(S\)-reductive group, and \(P,Q\)
two parabolic subgroups of \(G\).  The following conditions are
equivalent:
\begin{enumerate}[label=\textup{(\roman*)}]
\item \(P\cap Q\) is a parabolic subgroup of \(G\).
\item Étale-locally, \(P\cap Q\) contains a Borel subgroup of \(G\).
\item Étale-locally, \(P\cap Q\) contains a maximal torus of \(G\);
moreover, for every \(S'\to S\) and every maximal torus \(T\) of
\(G_{S'}\) contained in both \(P_{S'}\) and \(Q_{S'}\), the parabolic
subgroup opposite to \(P_{S'}\) relative to \(T\) is in transverse
position relative to \(Q_{S'}\).
\item There is an étale covering family \(\{S_i\to S\}\) and, for
each \(i\), a maximal torus \(T_i\) of \(G_{S_i}\), contained in both
\(P_{S_i}\) and \(Q_{S_i}\), such that the parabolic subgroup opposite
to \(P_{S_i}\) relative to \(T_i\) is in transverse position relative
to \(Q_{S_i}\).
\item There is an étale covering family \(\{S_i\to S\}\) and, for
each \(i\), a splitting \((T_i,M_i,R_i)\) of \(G_{S_i}\) such that
\(P_{S_i}\), respectively \(Q_{S_i}\), is the subgroup of type
\textup{(R)} of \(G_{S_i}\) containing \(T_i\) and defined by a set
of roots \(R_i^{(1)}\), respectively \(R_i^{(2)}\), with
\(R_i^{(1)}\cap R_i^{(2)}\) containing a system of positive roots
of \(R_i\).
\end{enumerate}
Moreover, when these conditions hold,
\[
 t(P\cap Q)=t(P)\cap t(Q)
\]
in the notation of \SGARefLink{sga3:xxvi:item:3:1}{3.2}.
\end{sgastatement}
\SGARefTarget{sga3:xxvi:prop:4:4:1}

The implications \textup{(v)}\(\Rightarrow\)\textup{(ii)} and
\textup{(iii)}\(\Rightarrow\)\textup{(iv)} are immediate.  Moreover,
\textup{(ii)}\(\Rightarrow\)\textup{(i)} follows from \SGARefLink{sga3:xxvi:rem:1:18}{1.18}.  To prove
\textup{(iv)}\(\Rightarrow\)\textup{(v)}, one may suppose that \(G\)
is split and that \(P\), respectively \(Q\), is defined by a set of
roots \(R_1\), respectively \(R_2\).  The parabolic subgroup opposite
to \(P\) is then defined by \(-R_1\), and the assertion reduces to
\SGARefLink{sga3:xxvi:lemma:4:2:2}{Lemma 4.2.2}.  The implication
\textup{(i)}\(\Rightarrow\)\textup{(iii)} is proved in the same way
after splitting.  Finally, the last assertion may be proved
étale-locally: one may suppose that \(P\cap Q\) contains a Borel
subgroup \(B\) of \(G\), and the assertion then reduces to \SGARefLink{sga3:xxvi:remark:3-7}{3.7}.

\begin{sgastatement}{Definition 4.4.2}
Two parabolic subgroups of \(G\) satisfying the equivalent conditions
\textup{(i)}--\textup{(v)} of \SGARefLink{sga3:xxvi:proposition:4-4-1}{4.4.1} are said to be in
\emph{osculating position}.
\end{sgastatement}
\SGARefTarget{sga3:xxvi:def:4:4:2}

\begin{sgastatement}{Corollary 4.4.3}
Let \(P,Q\) be parabolic subgroups in osculating position, and let
\(P',Q'\) be parabolic subgroups of \(G\) of the same types as \(P,Q\),
respectively.  The subgroups \(P',Q'\) are in osculating position if
and only if, étale-locally, the pair \((P',Q')\) is conjugate to the
pair \((P,Q)\).
\end{sgastatement}
\SGARefTarget{sga3:xxvi:cor:4:4:3}

It is enough to show that, if both \(P\) and \(P'\) are in osculating
position relative to \(Q\), then they are conjugate, étale-locally,
by a section of \(Q\).  Now \(P\cap Q\) and \(P'\cap Q\) are parabolic
subgroups of the same type contained in \(Q\).  By part \textup{(ii)}
of the lemma below, they are therefore conjugate étale-locally by a
section of \(Q\).  We may consequently suppose
\(P\cap Q=P'\cap Q\); part \textup{(i)} of that lemma then gives
\(P=P'\).

\begin{sgastatement}{Lemma 4.4.4}
Let \(P,P',Q\) be three parabolic subgroups of the \(S\)-reductive
group \(G\).
\begin{enumerate}[label=\textup{(\roman*)}]
\item The equality \(P=P'\) holds if and only if \(P\) and \(P'\) are
in osculating position and have the same type.
\item Suppose \(P\subset Q\) and \(P'\subset Q\).  If \(g\in G(S)\)
is such that \(\Int(g)P\) and \(P'\) are in osculating position, then
\(g\in Q(S)\).
\end{enumerate}
\end{sgastatement}
\SGARefTarget{sga3:xxvi:lemma:4:4:4}

Part \textup{(i)} follows immediately from the last assertion of
\SGARefLink{sga3:xxvi:proposition:4-4-1}{4.4.1}.  To prove \textup{(ii)}, observe that \(Q\) and \(\Int(g)Q\)
both contain \(P'\cap\Int(g)P\), and hence are in osculating position.
They coincide by \textup{(i)}, so that
\[
 g\in\Norm_G(Q)(S)=Q(S).
\]

Assertions \textup{(iii)} and \textup{(iv)} of \SGARefLink{sga3:xxvi:proposition:4-4-1}{Proposition 4.4.1} give
at once:

\begin{sgastatement}{Corollary 4.4.5}
Let \(P,P',Q\) be three parabolic subgroups of the \(S\)-reductive
group \(G\), containing the same maximal torus \(T\) of \(G\).
Suppose that \(P\) and \(P'\) are opposite relative to \(T\).  Then
\(Q\) is in osculating position relative to \(P\) if and only if it is
in transverse position relative to \(P'\).  Under these conditions,
\(P\cap Q\) is also in transverse position relative to \(P'\).
\end{sgastatement}
\SGARefTarget{sga3:xxvi:cor:4:4:5}

\begin{sgastatement}{Corollary 4.4.6}
Let \(P,Q\) be two parabolic subgroups of \(G\) containing the same
maximal torus \(T\).  The subgroups \(P,Q\) are in transverse position
if and only if there are parabolic subgroups \(P'\subset P\) and
\(Q'\subset Q\) of \(G\), opposite relative to \(T\).  One may even
choose
\[
 t(P')=t(P)\cap s(t(Q)).
\]
\end{sgastatement}
\SGARefTarget{sga3:xxvi:cor:4:4:6}

The condition is plainly sufficient, by \SGARefLink{sga3:xxvi:theorem:4-2-1}{4.2.1}\textup{\SGARefLink{sga3:xxvi:theorem:4-2-1:item:i}{(i)}} and
\SGARefLink{sga3:xxvi:item:4:3:1}{4.3.2}\textup{\SGARefLink{sga3:xxvi:item:4:3:1:item:iii}{(iii)}}.  Conversely, let \(P^-\), respectively \(Q^-\),
be the parabolic subgroup opposite to \(P\), respectively \(Q\),
relative to \(T\).  By \SGARefLink{sga3:xxvi:cor:4:4:5}{4.4.5}, \(P^-\cap Q\) is in transverse position
relative to both \(P\) and \(Q^-\), and therefore also relative to
\(P\cap Q^-\), by another application of \SGARefLink{sga3:xxvi:cor:4:4:5}{4.4.5}.  Moreover,
\[
\begin{aligned}
 t(P^-\cap Q)
 &=t(P^-)\cap t(Q)\\
 &=s(t(P))\cap s(t(Q^-))\\
 &=s\bigl(t(P)\cap t(Q^-)\bigr)\\
 &=s\bigl(t(P\cap Q^-)\bigr).
\end{aligned}
\]
It follows that
\[
 P^-\cap Q=P',\qquad P\cap Q^-=Q'
\]
are opposite, by \SGARefLink{sga3:xxvi:item:4:3:1}{4.3.2}\textup{\SGARefLink{sga3:xxvi:item:4:3:1:item:iii}{(iii)}}.  Since \(P'\cap Q'\supset T\),
they are indeed opposite relative to \(T\).
